\section{Approximation consequences}\label{sec:consequences}

Theorem~\ref{thm:ugc} turns established conditional hardness results
into NP-hardness results. The underlying conditional reductions and
approximation thresholds are due to the cited authors. We give the
additional arguments needed for the exact instance models stated below.
We also record two results from companion
manuscripts. For a maximization problem, ratio
\(r\) means returning value at least \(r\) times the optimum; for a
minimization problem, factor \(c\) means returning value at most \(c\)
times the optimum.
All factors and tolerances below are fixed independently of input size.
In each reduction using Theorem~\ref{thm:ugc}, they determine the
Unique Games errors before its alphabet is chosen. If
\(\mathrm P\ne\mathrm{NP}\), the resulting hardness statements exclude
deterministic polynomial-time approximations at the indicated strict
thresholds.

\subsection{Constraint satisfaction and Vertex Cover}

For \textsc{Max-Cut}, which maximizes the number of edges crossing a
bipartition, Khot, Kindler, Mossel, and O'Donnell's reduction gives
NP-hardness of approximation at every fixed ratio in
\((\alpha_{\mathrm{GW}},1)\), where the Goemans--Williamson constant is
\(\alpha_{\mathrm{GW}}\simeq0.87856\)~\cite[Theorem~1]{KKMO07}.
This matches the semidefinite approximation threshold of Goemans and
Williamson~\cite{GW95}. The reduction uses the Majority Is Stablest
theorem of Mossel, O'Donnell, and Oleszkiewicz~\cite{MOO10}.
For \textsc{Vertex Cover}, which minimizes the number of vertices
meeting every edge, Khot and Regev's reduction gives NP-hardness at
every fixed factor in \((1,2)\)~\cite{KhotRegev2008}.

A constraint satisfaction problem (CSP) is specified here by a finite
nonempty domain \(D\) and a finite family \(\Lambda\) of Boolean
predicates on \(D\), of bounded arity. A weighted \(\Lambda\)-instance
has a finite variable set and a finite collection of predicates from
\(\Lambda\) applied to tuples of variables, with nonnegative rational
weights summing to one. An assignment maps the variables to \(D\);
its value is the total satisfied weight, and \(\operatorname{OPT}\)
is the maximum of this value over all assignments.
Raghavendra's theorem characterizes the optimal approximation
for every fixed language of this form, up to arbitrarily small positive
error, by a basic semidefinite relaxation
\cite[Theorem~1.1 and Corollary~1.3]{Raghavendra08}.
Write \(\operatorname{LC}\) for the value of Raghavendra's LC relaxation,
as defined in~\cite[Section~4.5]{RaghavendraThesis09}.

The fixed-instance CSP theorem and the ordering reduction use the
following vertex-wise completeness guarantee.

Orient the constraints of a bipartite game from its left side to its
right side.  A left vertex is \emph{strongly satisfied} by a labeling
when every edge incident to it is satisfied.

\begin{lemma}[Simple games with completeness on left vertices]\label{lem:simple-left-vertex-completeness}
Fix \(0<\zeta,\nu<1\) and integers \(d,H\ge1\).  There are integers
\(h,R\ge1\), with \(h\ge H\) and \(d\mid h\), and a deterministic
polynomial-time reduction from \(\mathrm{3SAT}\) to nonempty explicit
unweighted simple bipartite Unique Games instances
\(\Phi_\varphi=(A,B,E)\) over \([R]\) with the following properties.
Every left vertex has degree \(h\).  On a satisfiable input \(\varphi\),
one labeling strongly satisfies at least \((1-\zeta)|A|\) left vertices.
On an unsatisfiable input, \(\val(\Phi_\varphi)<\nu\).
The degree \(h\) and the two errors used in
Theorem~\ref{thm:ugc} may be chosen from \(\zeta,\nu,d,H\) before that
theorem supplies the alphabet \(R\).
\end{lemma}

Appendix~\ref{subsec:left-vertex-completeness-proof} proves this formulation, with
the errors and left degree chosen before the alphabet.

\begin{corollary}[Raghavendra's consequence]\label{cor:raghavendra-consequence}
Fix a finite domain \(D\), a finite bounded-arity Boolean predicate
family \(\Lambda\) on \(D\), and a weighted \(\Lambda\)-instance \(J\)
in the rational model above.  Suppose
\(A=\operatorname{OPT}(J)<B=\operatorname{LC}(J)\).
For every fixed \(0<\eta<(B-A)/2\), it is NP-hard to distinguish
\(\Lambda\)-instances \(I\) satisfying
\[
 \operatorname{OPT}(I)\ge B-\eta
 \qquad\text{from those satisfying}\qquad
 \operatorname{OPT}(I)\le A+\eta.
\]
\end{corollary}
The proof in Appendix~\ref{subsec:raghavendra-proof} applies the
fixed-instance theorem with tolerance $\eta/4$ and changes its fixed
finite verifier law by total variation less than $\eta/2$.  The latter
change yields rational weights and changes every assignment value by
less than $\eta/2$.


\subsection{Ordering constraints}

An ordering CSP uses a global order of the input variables, rather than
an assignment from a fixed finite domain. Fix an integer \(k\ge2\) and
a nonempty proper subset \(\Pi\subset S_k\), where \(S_k\) is the set
of permutations of \(\{1,\ldots,k\}\). In the finite
probability-weighted model of Guruswami, H\aa stad, Manokaran,
Raghavendra, and Charikar~\cite[Definition~8.2]{GuruswamiEtAl2011OrderingCSP},
we use an explicitly listed finite variable set \(V\) and a finite list of
constraint occurrences. Their nonnegative rational weights are encoded
by binary numerators and denominators and sum to one. Every occurrence is an
ordered \(k\)-tuple \((v_1,\ldots,v_k)\) of distinct variables,
satisfied by a linear order \(\prec\) when
\[
 v_{\pi(1)}\prec\cdots\prec v_{\pi(k)}
 \quad\text{for some }\pi\in\Pi.
\]
Write \(\OPT_\Pi(I)\) for the maximum satisfied weight. A uniform
random order satisfies expected weight \(\rho_\Pi=|\Pi|/k!\).
Since \(\OPT_\Pi(I)\le1\), it is a \(\rho_\Pi\)-approximation.
For fixed \(k\), conditional expectation gives a deterministic
polynomial-time algorithm with this guarantee. The end of
Appendix~\ref{subsec:ordering-transfer-proof} gives the exact
conditional probabilities and the bit-length bound.
The size of \(V\) is part
of the input, so this is not a CSP over the fixed finite domain used
in the preceding corollary.

\begin{proposition}[Finite ordering transfer]
\label{prop:finite-ordering-transfer}
Fix \(k\ge2\) and \(\varnothing\ne\Pi\subsetneq S_k\), put
\(\rho_\Pi=|\Pi|/k!\), and fix
\(0<\delta<(1-\rho_\Pi)/2\).  There are integers
\(q,H\ge1\) and positive rationals \(\zeta,\nu<1\), depending only
on \(k,\Pi,\delta\), with the following property for every fixed
integer \(R\ge1\).

Let \(\Phi=(A,B,E)\) be a nonempty explicit unweighted simple
bipartite Unique Games instance over \([R]\), oriented from \(A\) to
\(B\), whose common left degree \(h\) satisfies \(h\ge H\) and
\(q\mid h\).  There is a deterministic transformation, polynomial in
the explicit size of \(\Phi\) for fixed \(R\), producing an instance \(I\) given by an explicit
list of \(\Pi\)-ordering constraints on distinct variables within
each constraint, with nonnegative rational weights, each encoded by
a binary numerator and a positive binary denominator, summing to one.
If some labeling strongly satisfies at least
\((1-\zeta)|A|\) left vertices, then \(\OPT_\Pi(I)\ge1-\delta\).
If \(\val(\Phi)\le\nu\), then
\(\OPT_\Pi(I)\le\rho_\Pi+\delta\).
The constants \(q,H,\zeta,\nu\) are independent of both \(R\) and
\(h\).
\end{proposition}

Appendix~\ref{subsec:ordering-transfer-proof} proves the proposition.
The degree conditions can be supplied by Lemma~\ref{lem:simple-left-vertex-completeness}
before its alphabet is chosen.

\begin{corollary}[Ordering-CSP consequence]\label{cor:ordering-csp}
Fix \(k\ge2\) and \(\varnothing\ne\Pi\subsetneq S_k\), and put
\(\rho_\Pi=|\Pi|/k!\). For every fixed
\(0<\varepsilon<(1-\rho_\Pi)/2\), it is NP-hard to distinguish finite
weighted \(\Pi\)-ordering-CSP instances \(I\) satisfying
\[
 \OPT_\Pi(I)\ge1-\varepsilon
 \qquad\text{from those satisfying}\qquad
 \OPT_\Pi(I)\le\rho_\Pi+\varepsilon.
\]
\end{corollary}
Consequently, approximation at every fixed ratio
\(r\in(\rho_\Pi,1]\) is NP-hard. Together with the random-ordering
guarantee, this gives the optimal threshold \(\rho_\Pi\).
\begin{proof}
Apply Proposition~\ref{prop:finite-ordering-transfer} with
\(\delta=\varepsilon\). It uses the local gap and analytic estimates of
Guruswami et al.~\cite{GuruswamiEtAl2011OrderingCSP} and supplies
\(q,H,\zeta,\nu\), depending only on \(k,\Pi,\varepsilon\), before any
alphabet is chosen. Apply Lemma~\ref{lem:simple-left-vertex-completeness} with these
\(\zeta,\nu\), with \(d=q\), and with the same \(H\). The degree and errors
are fixed before the alphabet is chosen. Using
Theorem~\ref{thm:ugc}, it maps each \(\mathrm{3SAT}\) formula \(\varphi\)
to a simple game with completeness on left vertices \(\Phi_\varphi=(A,B,E)\) over a fixed
alphabet \([R]\) whose
common left degree \(h\) satisfies \(h\ge H\) and \(q\mid h\). If
\(\varphi\) is satisfiable, one labeling strongly satisfies at least
\((1-\zeta)|A|\) left vertices; otherwise \(\val(\Phi_\varphi)<\nu\).
These are the hypotheses of
Proposition~\ref{prop:finite-ordering-transfer} for this fixed \(R\), so
composing the two reductions gives \(\OPT_\Pi(I)\ge1-\varepsilon\) in the
first case and \(\OPT_\Pi(I)\le\rho_\Pi+\varepsilon\) in the second. This
proves the displayed gap for the original predicate \(\Pi\).

For a fixed \(r\in(\rho_\Pi,1]\), choose the positive error still
smaller so that \(\varepsilon<(r-\rho_\Pi)/(1+r)\). Then
\(r(1-\varepsilon)>\rho_\Pi+\varepsilon\), so an
\(r\)-approximation would distinguish the two cases.
\end{proof}

For \textsc{Maximum Acyclic Subgraph} in the finite weighted
directed-graph model, take \(k=2\) and \(\Pi=\{12\}\) in one-line
notation. A constraint \((u,v)\) asks that the edge point forward;
forward edges form an acyclic subgraph, and every acyclic subgraph has
a topological order. Its threshold is therefore \(1/2\).
For \textsc{Betweenness}, the constraint \((a,b,c)\) asks that \(b\)
lie between \(a\) and \(c\), so \(\Pi=\{123,321\}\subset S_3\) and
the threshold is \(2/3!=1/3\)
\cite[Theorem~1.1 and Section~1.2]{GuruswamiEtAl2011OrderingCSP}.
These finite weighted consequences have near-perfect completeness
with fixed positive error.

\subsection{Cut, deletion, and clustering consequences}

For the cut problems of Chawla, Krauthgamer, Kumar, Rabani, and
Sivakumar~\cite[Section~1]{CKKRS2006}, let \(G=(V,E)\) be a finite
undirected graph with positive rational edge costs \(c_e\), and let
\(\mathcal D=\{\{s_i,t_i\}:1\le i\le k\}\) contain \(k\ge1\)
distinct demand pairs with \(s_i\ne t_i\). For \(M\subseteq E\), write
\(c(M)=\sum_{e\in M}c_e\), and let \(q_{\mathcal D}(M)\) count
the demand pairs disconnected in \(G-M\).
Weighted \textsc{Multicut} minimizes \(c(M)\) subject to
\(q_{\mathcal D}(M)=k\). Their nonuniform \textsc{Sparsest Cut}
search problem minimizes \(c(M)/q_{\mathcal D}(M)\) over cutsets
\(M\) with \(q_{\mathcal D}(M)>0\). Only the selected pairs carry
unit demand; they need not be all vertex pairs.

For weighted Min-\(2\mathrm{CNF}^{\equiv}\) Deletion, the input
consists of clauses \((\ell\equiv\ell')\) with positive rational
weights, where \(\ell,\ell'\) are Boolean literals, and an assignment
pays the total weight of false clauses. The allowed clauses therefore
express equality or disequality of variables.

In the arbitrary weighted-graph form of \textsc{Correlation Clustering},
each present edge has a \(+\) or \(-\) label and a nonnegative rational
weight. A partition into any number of clusters pays for \(+\) edges
whose endpoints lie in different clusters and \(-\) edges whose
endpoints lie in the same cluster. Missing edges carry no penalty.
This is the minimum-disagreements objective of
\cite[Section~2.1]{DemaineEtAl2006Clustering}.

\begin{corollary}[Cut, deletion, and clustering hardness]
\label{cor:cut-clustering}
For every fixed real constant \(C>1\), it is NP-hard to approximate
weighted \textup{\textsc{Multicut}}, nonuniform \textup{\textsc{Sparsest Cut}}, weighted
Min-\(2\mathrm{CNF}^{\equiv}\) Deletion, and weighted
\textup{\textsc{Correlation Clustering}}, in the models above, within factor \(C\).
For \textup{\textsc{Sparsest Cut}}, this is hardness of the search task under
polynomial-time Cook reductions: the algorithm must output a cutset,
not only an estimate of the optimum.
\end{corollary}
\begin{proof}
Corollary~1.3 of \cite{CKKRS2006} gives the first three conclusions
under its ordinary edge-value Unique Games hypothesis.
For a requested factor \(C\), use its hardness statement at any fixed
larger factor \(L\): a \(C\)-approximation would also be an
\(L\)-approximation.
To supply that hypothesis, orient the bipartite constraints of
Theorem~\ref{thm:ugc} left-to-right, renumber its fixed alphabet as
\([d]\) with \(d=2^s\ge2\), and give its edges equal weights summing
to one. The unique 2-prover games of \cite[Section~1.2]{CKKRS2006} have the
same number of questions on each side, so pad the smaller vertex side
with isolated questions if needed; this preserves both value and
polynomial size. The search-version
qualification for \textsc{Sparsest Cut} is in Section~1.3 of the
cited source.

For \textsc{Correlation Clustering}, the hard \textsc{Multicut}
instances in~\cite[Section~2.1]{CKKRS2006} have distinct antipodal
demand pairs inside \(d\)-cubes, while the edges of the Multicut graph
are cube edges or join different cubes. Lemma~1.5 of that source, which makes the game regular, does not
change \(d\), so \(d\ge2\) there. Antipodal vertices of a \(d\)-cube
differ in all \(d\ge2\) coordinates, whereas a cube edge joins vertices
that differ in exactly one; edges between cubes cannot join
a same-cube demand pair either. Thus the demand pairs are nonedges. The transformation of
Demaine, Emanuel, Fiat, and Immorlica
\cite[Theorem~4.7 and Corollary~4.8]{DemaineEtAl2006Clustering}
therefore creates no opposite-signed parallel edges. It preserves the
optimum and converts every clustering into a multicut of no greater
cost, so the same fixed-factor hardness transfers.
To see the conversion explicitly, put \(W_+=\sum_{e\in E}c_e\).
Keep the original edges as positive edges with their costs, and give
each demand pair a negative edge of weight \(W_++1\).
If a clustering separates every demand pair, its crossing positive
edges form a multicut of exactly its disagreement cost.  If it leaves
a demand pair together, its cost exceeds \(W_+\); returning all
original edges is then a multicut of smaller cost.  Conversely, the
connected components left by a multicut form a clustering of cost at
most that multicut's cost.  Taking minima in both directions proves
equality of optima.  The sums and the added weights have polynomial
binary length, so this also preserves the rational input model.
\end{proof}

The separate manuscript
\cite[Theorem~1.1]{OpenAIUniformSparsest2026} proves
NP-hardness at every fixed factor \(C>1\) for the bipartition objective with
nonnegative rational capacities and exactly unit demand between every
pair of distinct vertices.

\paragraph{Kernel clustering.}
A separate companion application concerns the following maximization
problem. For \emph{identity-target kernel clustering}, fix an integer \(k\ge3\).
An input is an explicitly represented rational symmetric matrix
\(A=(a_{pq})_{p,q=1}^N\succeq0\) with \(A\mathbf1=0\), where
\(\mathbf1\) is the all-ones vector.
An assignment \(\sigma:\{1,\ldots,N\}\to\{1,\ldots,k\}\) has value
\[
 \val_A(\sigma)=\sum_{p,q=1}^N
 a_{pq}\mathbf1_{\{\sigma(p)=\sigma(q)\}},
 \qquad
 \OPT_k(A)=\max_\sigma\val_A(\sigma).
\]
The sum includes ordered pairs and the diagonal; empty clusters are
allowed, and there is no normalization by cluster size. A loss factor
\(\alpha\ge1\) means returning value at least \(\OPT_k(A)/\alpha\).
The Gaussian propeller companion uses Theorem~\ref{thm:ugc} to prove
NP-hardness of achieving every fixed loss factor
\[
 1\le\alpha<\alpha_k,
 \qquad \alpha_k=\frac{8\pi}{9}\left(1-\frac1k\right),
\]
for each fixed \(k\ge3\)
\cite[Theorem~6.1]{OpenAIGaussianPropeller2026}.
This is strict hardness below the rounding factor
\(\alpha_k\).
